% Осень26!A3; нумерация внутри столбца без пустых строк.
% Источник: лекция 1, §1 (строка 29)
% Источник: лекция 1, §2 (строка 66)
% Источник: лекция 1, §3 (строка 106)
% Источник: лекция 1, §4 (строка 122)
% Источник: лекция 1, §5 (строка 136)
\ticket{01}{Независимость частных производных от порядка дифференцирования}
\label{ticket:01}

\subsection*{Вторые и высшие частные производные}

\subsubsection*{Напоминание: дифференцируемость и частные производные}

Дифференцируемость \(f\) в точке \(a\) --- это возможность \textbf{линейного приближения}: \(f\) = константа + линейная функция + \(o(\cdot)\)-малое в окрестности \(a\). Из дифференцируемости следует существование частных производных; в обратную сторону «хитрее»: частные производные могут существовать, а функция не быть дифференцируемой (частные производные --- «приличное поведение при движении в точку параллельно осям», при произвольном движении «что-то сбивается»). Теорема прошлого семестра: если частные производные существуют в окрестности и \textbf{непрерывны}, то из этого следует дифференцируемость.

Теперь «освоились с однократной операцией дифференцирования» --- добавляем \textbf{многократное}.

\subsubsection*{Определение второй частной производной}

\begin{definition}

Пусть \(f:E\subset\mathbb{R}^m\to\mathbb{R}\) (числовая функция, «\(f\) маленькое» --- не отображение), \(a\in\operatorname{Int}E\). Фиксируем номер координаты \(k\in\overline{1,m}\). Пусть существует окрестность \(U(a)\subset E\), в каждой точке которой определена частная производная \(\dfrac{\partial f}{\partial x_k}\), т.е. определена функция \[g:=\frac{\partial f}{\partial x_k}:U(a)\to\mathbb{R}.\] Если существует частная производная \(\dfrac{\partial g}{\partial x_i}(a)\), то она называется \textbf{второй частной производной} \(f\) по переменным \(x_k\) и \(x_i\) (в этом порядке): \[\frac{\partial^2 f}{\partial x_i\,\partial x_k}(a):=\frac{\partial}{\partial x_i}\!\left(\frac{\partial f}{\partial x_k}\right)(a).\]

\end{definition}

\textbf{О порядке дифференцирований.} В этом определении порядок зафиксирован: \textbf{первым} дифференцируем по \(x_k\), \textbf{затем} по \(x_i\). Порядок «поддерживается выписыванием дифференциальчиков справа налево»: в записи \(\frac{\partial^2 f}{\partial x_i\,\partial x_k}\) правый сомножитель \(\partial x_k\) отвечает первому дифференцированию. Логика обозначения «наивная»: оператор \(\frac{d}{dx_i}\) применяется к результату оператора \(\frac{d}{dx_k}\), и дроби «перемножаются»: \(\frac{d}{dx_i}\cdot\frac{d}{dx_k}=\frac{d^2}{dx_i\,dx_k}\).

\textbf{Компактные обозначения} (старомодное и короткое): \(f''_{x_i x_k}(a)\), где переменные перечисляются в том же порядке (справа налево по времени дифференцирования); если ясно из контекста, можно писать только номера: \(f''_{ik}\). Для старших порядков --- аналогично: \(f^{(r)}_{x_{i_1}\ldots x_{i_r}}\) или \(f^{(r)}_{i_1\ldots i_r}\), а также \(f^{(r)}\) с номером \(r\) = порядок производной.

\subsubsection*{Высшие частные производные}

\begin{definition}

\emph{рекурсивное.} Пусть в некоторой окрестности точки \(a\) определена функция \[g=\frac{\partial^{r-1} f}{\partial x_{i_2}\cdots\partial x_{i_r}}.\] Тогда \[\frac{\partial^{r} f}{\partial x_{i_1}\partial x_{i_2}\cdots\partial x_{i_r}}(a):=\frac{\partial}{\partial x_{i_1}}\!\left(\frac{\partial^{r-1}f}{\partial x_{i_2}\cdots\partial x_{i_r}}\right)(a).\] Чтобы определение сработало, нужно, чтобы предыдущая \((r-1)\)-я производная была задана \textbf{в целой окрестности} точки \(a\) (дифференцирование --- всегда «операция в точке», и лишь убедившись, что оно выполняется всюду в окрестности, мы получаем функцию, которую можно дифференцировать дальше).

\end{definition}

\textbf{Замечание.} Следующая теорема (§2) скажет, что в «приличной ситуации» неважно, в каком порядке дифференцировать --- так что в будущем порядок в обозначениях перестаёт быть существенным; повторяющиеся переменные можно группировать: вместо «слова» \(x_1x_1x_1\) писать \(x_1^3\) (это не арифметическое выражение, а перечень букв, по которым происходит дифференцирование, --- с кратностями).

\subsubsection*{Пример}

\[f(x,y,z)=x^{z}+\frac{z}{y},\qquad \frac{\partial^2 f}{\partial x\,\partial y}=?\]

Читаем порядок справа налево: \textbf{сначала} дифференцируем по \(y\), потом по \(x\). По \(y\): слагаемое \(x^z\) от \(y\) не зависит, \(\dfrac{\partial}{\partial y}\dfrac{z}{y}=-\dfrac{z}{y^2}\) --- в результате \textbf{иксов нет}, поэтому после дифференцирования результата по \(x\) получаем \[\frac{\partial^2 f}{\partial x\,\partial y}=0 .\] (\textbf{Требует сверки.} Функция восстановлена из шумного транскрипта «X в степени Z + Z / Y»; если там было \(xz+z/y\), ответ и ход рассуждения не меняются.)

\subsection*{Теорема о независимости частных производных от порядка дифференцирования (двумерный случай)}

Лектор, как и прежде для громоздких утверждений, доказывает случай двух переменных («в нём всё видно, а увеличение размерности делает всё громоздким, но не сложнее»).

\begin{theorem}

Пусть \(f:E\subset\mathbb{R}^2\to\mathbb{R}\), \((x_0,y_0)\in\operatorname{Int}E\). Пусть существует окрестность \(U(x_0,y_0)\), в которой определены (как функции) обе смешанные производные \(f''_{xy}\) и \(f''_{yx}\), и обе они \textbf{непрерывны в точке} \((x_0,y_0)\). Тогда \[f''_{xy}(x_0,y_0)=f''_{yx}(x_0,y_0).\]

\end{theorem}

\textbf{Замечания к формулировке.}

\begin{itemize}
\tightlist
\item
  Здесь одноточечный вариант: непрерывность требуется только в центральной точке, и совпадение получается в ней же. Если обе производные непрерывны \textbf{во всей окрестности}, то они совпадают \textbf{во всей окрестности} (глобальный вариант). (Лектор в начале колебался, где требовать непрерывность --- «в центральной точке\ldots{} если будет не хватать, поправим на окрестность»; принята формулировка справочника.)
\item
  Напоминание о соглашении: \(f''_{yx}\) --- сначала дифференцируем по \(x\) (правый нижний индекс), затем по \(y\); очередное дифференцирование \textbf{дописывается слева}.
\end{itemize}

\textbf{Идея доказательства.} Рассматривается «вторая разность» \(\Delta^2(h,k)\) --- знакочередующаяся сумма значений \(f\) в вершинах прямоугольника. Она вычисляется двумя способами: дважды применяется одномерная теорема Лагранжа --- сначала по переменной \(x\) (с фиксацией \(k\)), потом по \(y\); затем те же два применения в противоположном порядке. Первый способ даёт \(\Delta^2=f''_{yx}\,hk\) в средней точке одного вида, второй --- \(\Delta^2=f''_{xy}\,hk\) в средней точке другого вида. После деления на \(hk\) и предельного перехода \((h,k)\to(0,0)\) непрерывность смешанных производных «протаскивает» оба выражения в точку \((x_0,y_0)\), где они и совпадают.

\begin{proof}

Возьмём приращения \(h,k\) столь малыми, чтобы весь прямоугольник с вершинами \((x_0,y_0)\), \((x_0+h,y_0)\), \((x_0,y_0+k)\), \((x_0+h,y_0+k)\) лежал в указанной окрестности (для круглой окрестности радиуса \(\rho\) достаточно, например, \(|h|,|k|\le\rho/\sqrt2\)). Определим \textbf{вторую разность} \[\Delta^2(h,k):=f(x_0+h,\,y_0+k)-f(x_0+h,\,y_0)-f(x_0,\,y_0+k)+f(x_0,\,y_0)\] --- значения «с плюсиками» берутся по одной диагонали прямоугольника, «с минусиками» --- по другой.

\textbf{Эксперимент 1 (фиксируем \(k\), работаем по \(h\)).} Пусть \(\alpha(h):=\Delta^2(h,k)\) при фиксированном \(k\). Тогда \(\alpha(0)=0\) (слагаемые попарно сокращаются), и \[\alpha(h)=\alpha(h)-\alpha(0).\] Функция \(u\mapsto f(u,y_0+k)-f(u,y_0)\) дифференцируема по \(u\) (по первой переменной \(f\) дифференцируема в всей окрестности --- там даже вторые производные определены), поэтому по одномерной \textbf{теореме Лагранжа} \[\alpha(h)=\alpha'(\tilde h)\,h,\qquad \tilde h\ \text{между }0\text{ и }h,\] \[\alpha'(\tilde h)=f'_x(x_0+\tilde h,\,y_0+k)-f'_x(x_0+\tilde h,\,y_0).\] Важно: \textbf{одна и та же} средняя точка \(\tilde h\) стоит в обоих слагаемых --- производная \(\alpha'\) есть разность значений, и \(\tilde h\) подставляется туда одновременно. Теперь в этой разности фиксирована первая координата \(x_0+\tilde h\), а вторая принимает два значения \(y_0\) и \(y_0+k\) --- снова теорема Лагранжа, теперь по второй координате: \[\alpha'(\tilde h)=f''_{yx}(x_0+\tilde h,\,y_0+\tilde k)\cdot k,\qquad \tilde k\ \text{между }0\text{ и }k\] (производная по \(y\) от \(f'_x\) --- это \(f''_{yx}\): к \(f'_x\) дописываем дифференцирование по \(y\) слева). Итого \[\Delta^2(h,k)=f''_{yx}(x_0+\tilde h,\,y_0+\tilde k)\,h\,k.\tag{1}\]

\textbf{Эксперимент 2 (фиксируем \(h\), работаем по \(k\)).} Пусть \(\beta(k):=\Delta^2(h,k)\) при фиксированном \(h\). Тот же ход выкладок в другом порядке даёт \[\Delta^2(h,k)=f''_{xy}(x_0+\hat h,\,y_0+\hat k)\,h\,k,\tag{2}\] со «своими» средними точками \(\hat h,\hat k\) (между \(0\) и \(h\), соответственно \(0\) и \(k\)). «Гипнотизёр» с фиксацией --- лишь способ следить за выкладкой: на самом деле мы взяли одно и то же выражение \(\Delta^2(h,k)\) и применили теорему Лагранжа дважды, по-разному группируя.

\textbf{Предел.} При \(h\ne0\), \(k\ne0\) приравниваем (1) и (2) и делим на \(hk\): \[f''_{yx}(x_0+\tilde h,\,y_0+\tilde k)=f''_{xy}(x_0+\hat h,\,y_0+\hat k).\] Устремим \((h,k)\to(0,0)\); тогда \(\tilde h,\hat h\to0\) и \(\tilde k,\hat k\to0\) (средние точки зажаты между \(0\) и \(h\), соответственно \(0\) и \(k\)), т.е. обе точки, в которых вычислены производные, стремятся к \((x_0,y_0)\). По непрерывности \(f''_{yx}\) и \(f''_{xy}\) в точке \((x_0,y_0)\): \[f''_{yx}(x_0,y_0)=f''_{xy}(x_0,y_0).\]

\end{proof}

\textbf{Пояснение лектора о строгости предельного перехода.} Равенство (1)=(2) после деления на \(hk\) известно при всех малых \(h,k\ne0\); предел \(\Delta^2(h,k)/(hk)\) при \((h,k)\to(0,0)\) --- это \textbf{двойной предел} по проколотой окрестности нуля, из которой выколоты прямые \(h=0\) и \(k=0\) («выколотый крест»: знаменатель не должен обращаться в ноль). Формально можно рассмотреть функцию \((h,k)\mapsto\Delta^2(h,k)/(hk)\), сузив её область определения на окрестность с выколотым крестом, --- тогда всё корректно по общему определению предела. А так как непрерывность даёт обычный \textbf{метрический} предел (по норме), а из существования метрического предела следует существование и равенство двойного (повторного) --- рассуждение корректно; подробности см. §3.

\subsection*{Техническое отступление: двойной предел и «выколотый крест»}

Лектор счёл нужным явно проговорить, о каком пределе идёт речь (вопрос из аудитории: «откуда вы знаете слово „повторный предел``?»).

\textbf{Метрический предел} (на нём основана непрерывность в \(\mathbb{R}^m\)): \(F(h,k)\to L\) при \((h,k)\to(0,0)\) означает \[\forall\varepsilon>0\ \exists\delta>0\ \forall(h,k):\ 0<|(h,k)|<\delta\ \Rightarrow\ |F(h,k)-L|<\varepsilon,\] где \(|(h,k)|=\sqrt{h^2+k^2}\) --- расстояние до нуля; сама точка \((0,0)\) (центр) брать запрещено.

\textbf{Двойной (повторный) предел} здесь понимается так: \(h\) и \(k\) --- две \textbf{независимые} переменные, каждая пробегает проколотую окрестность нуля: \[\forall\varepsilon>0\ \exists\delta>0\ \forall h\,\bigl(0<|h|<\delta\bigr)\ \forall k\,\bigl(0<|k|<\delta\bigr)\ \Rightarrow\ |F(h,k)-L|<\varepsilon,\] т.е. \((h,k)\) лежит в квадратной окрестности размера \(\delta\), из которой выброшен «крест» --- точки с \(h=0\) или \(k=0\). Это в точности тот же предел, что и выше, но «природа предела подменена»: определение через две независимые одномерные проколотые окрестности.

\textbf{Как устроено применение в §2.} У функции \((h,k)\mapsto \Delta^2(h,k)/(hk)\) область определения --- окрестность с выколотым крестом (вне креста знаменатель ненулевой). На этой суженной области предел существует (он метрический, раз производные непрерывны --- «непрерывность оперируется более сильным, метрическим пределом; раз метрический существует, то и двойной существует» --- обсуждали в прошлом семестре). Больше того: после сокращения на \(hk\) обе части равенства (1)=(2) определены и корректно продолжаются и на кресте, так что итоговое равенство непрерывных функций «не ранит» и включение точек креста.

\subsection*{\texorpdfstring{Классы гладкости \(C^r(E)\)}{Классы гладкости C\^{}r(E)}}

\begin{definition}

Пусть \(E\) --- открытое подмножество \(\mathbb{R}^m\), \(r\in\mathbb{N}\). Класс \(C^r(E)\) --- множество всех функций \(f:E\to\mathbb{R}\), у которых \textbf{существуют все частные производные всех порядков до \(r\) включительно во всех точках \(E\)} и \textbf{все они непрерывны на \(E\)}.

\end{definition}

\textbf{Замечания.}

\begin{itemize}
\tightlist
\item
  Производные порядков \(\le r-1\) автоматически непрерывны: они дифференцируемы (существуют производные следующего порядка), а дифференцируемость даёт непрерывность. Содержательное требование --- непрерывность производных ровно \(r\)-го порядка.
\item
  \(C(E)\) (без индекса) --- множество всех непрерывных функций на \(E\); это то же, что \(C^0(E)\).
\item
  Возникает цепочка вложений \[C(E)\supset C^1(E)\supset C^2(E)\supset\cdots\supset C^r(E)\supset\cdots,\qquad \bigcap_{r=1}^{\infty}C^r(E)=:C^\infty(E)\] --- класс \textbf{бесконечно дифференцируемых} функций.
\item
  Все включения \textbf{собственные}. Примеры: \(f(x)=|x_1|\) (прочие переменные игнорируются, \(0\in E\) --- внутренняя точка) непрерывна, но \(\partial f/\partial x_1\) не существует в точках с \(x_1=0\) --- значит \(f\in C(E)\setminus C^1(E)\); примеры из одномерной ситуации «поднимаются в \(\mathbb{R}^m\)». Функция \(f(x)=x\) --- бесконечно дифференцируемая (при дальнейших дифференцированиях «ноль получается, но всё равно она бесконечно дифференцируется»), так что пересечение классов непусто --- «теоремы Кантора про убывающие множества нет, общая точка может быть, и вот она».
\end{itemize}

\subsection*{Общая теорема о независимости частных производных от порядка дифференцирования}

\begin{theorem}

\emph{общий случай.} Пусть \(E\subset\mathbb{R}^m\) открыто, \(f\in C^r(E)\), \(r\ge 2\) (при \(r=1\) «менять порядок» ещё нечего). Пусть \((i_1,\dots,i_r)\) и \((j_1,\dots,j_r)\) --- два набора номеров координат (каждый номер из \(\overline{1,m}\)), отличающиеся \textbf{перестановкой}. Тогда \[\frac{\partial^r f}{\partial x_{i_1}\dots\partial x_{i_r}}=\frac{\partial^r f}{\partial x_{j_1}\dots\partial x_{j_r}}\qquad\text{на всём множестве } E\] (одноточечный вариант теоремы §2 в силу открытости \(E\) и непрерывности автоматически становится глобальным).

\end{theorem}

\textbf{Идея доказательства (лекция --- «замечания в сторону общей теоремы»; полная схема --- справочник из исходного конспекта §1.2).} Всякая перестановка набора порождается не просто транспозициями, а \textbf{транспозициями соседних элементов} («перейти от первой перестановки к другой --- сделав несколько операций „поменять местами два соседних индекса``»). Поэтому достаточно уметь менять местами два \textbf{соседних} дифференцирования. Если «внешние» дифференцирования (до и после выбранной пары) пока «не сыграли» --- отложить их: обозначить уже вычисленную частную производную порядка \(r-2\) новой буквой, скажем \(F\), и применить к \(F\) \textbf{двумерную теорему §2}: \(F''_{x_p x_q}=F''_{x_q x_p}\) (обе смешанные производные \(F\) существуют и непрерывны, поскольку \(f\in C^r\)). Раз перестановочна одна пара соседних индексов --- перестановочны и остальные композиции таких перестановок: «ключевое --- аналитическая причина (двумерный случай), дальше чистая комбинаторика».
